\fancyhead[L]{\small FUNDAMENTALGLEICHUNGEN.}
\SourceLeaf{501}\pdfbookmark[0]{Fundamentalgleichungen für die Bewegung schwerer Körper auf der rotierenden Erde.}{rotating-earth-title}
\begin{center}
{\large\bfseries FUNDAMENTALGLEICHUNGEN\par\medskip FÜR DIE BEWEGUNG SCHWERER KÖRPER\par\medskip AUF DER ROTIERENDEN ERDE.\par}
\bigskip
\textsc{Benzenberg}. Versuche über das Gesetz des Falls. 1804.
\par\bigskip
{\large\itshape Brief von Gauss an Benzenberg.\par}
\end{center}
\begin{flushright}Braunschweig 1803. 2. Februar.\end{flushright}
—\quad—\quad—\quad In der Theorie unseres Freundes \textsc{Olbers} steckt eine Voraussetzung, die mir unzulässig scheint. Nämlich: \textit{dass der Körper während des Falls in einer Ebene bleibt.} Das darf man meiner Meinung nach aber \textit{nicht} voraussetzen, wenn man den Luftwiderstand berücksichtigt. Den muss man hier \textit{unbedingt} berücksichtigen, weil die erschlossene Abweichung nach Süden allein darauf beruht. Eine einfache Überlegung zeigt Folgendes: Die Ebene $(\mathrm{A})$, in der sich der Körper anfangs zu bewegen beginnt, geht durch den Mittelpunkt der Erde (oder allgemeiner durch den Mittelpunkt der Anziehung) und steht senkrecht auf der Ebene $(\mathrm{B})$, in der zu Beginn des Falls der Meridian des Beobachtungsorts lag. Man sieht aber leicht, dass die Luftteilchen an allen Stellen schräg durch die Ebene $\mathrm{A}$ gehen, außer auf der Geraden, auf der $\mathrm{A}$ von $\mathrm{B}$ geschnitten wird. Die Luft wirkt dem Körper daher in einer Richtung außerhalb dieser Ebene $\mathrm{A}$ entgegen und treibt ihn nach Norden aus ihr heraus. Mir schien, dass dieser Effekt gerade groß genug wäre, um die aus der Verzögerung des Falls erschlossene Abweichung nach Süden aufzuheben.

Ihr letzter Brief veranlasste mich, darüber erneut nachzudenken. In einer freien halben Stunde betrachtete ich die Sache auf eine ganz andere Weise und leitete aus den ersten Grundsätzen der Dynamik die analytischen Gleichungen her,
\SourceLeaf{502}die die relative Bewegung des Körpers gegenüber der bewegten Erdoberfläche beschreiben. Dabei fand ich zu meiner Verwunderung

1) die Abweichung nach Süden wieder $0$ oder ganz unmerklich.

2) die Abweichung nach Osten nur $\tfrac23$ von dem, was Dr.\ \textsc{Olbers} gefunden hat. In der Schreibweise von Dr.\ \textsc{Olbers} ist sie, wenn man den Luftwiderstand vernachlässigt,
\[
=\frac{\tfrac43\pi\cos\psi.at}{86164}
\]
oder, wenn man ihn berücksichtigt, mit einer hier ausreichenden Näherung
\[
=\frac{\tfrac43\pi\cos\psi.t}{86164}(\tfrac32a'-\tfrac12a)
\]
wobei $a$ die Höhe ist, die der Körper in der Zeit $t$ im Vakuum durchfallen würde, also $=\tfrac12g'tt$ [wobei ferner $\psi$ die Polhöhe des Beobachtungsorts und $a'$ die tatsächliche Fallhöhe bezeichnet].

Setze ich demnach für Ihre Versuche die Pendellänge für Hamburg $=440.75$ Linien (woraus sich $g'$ fast so ergibt, wie Dr.\ \textsc{Olbers} es annimmt), so finde ich \textit{die Abweichung nach Osten} zu $3.951$ \textit{Pariser Linien}. Das stimmte sehr genau mit Ihren Versuchen überein, — während die Abweichung nach Süden zu meinen Ergebnissen nicht passt.

\begin{center}\resizebox{!}{58mm}{\begin{tikzpicture}[x=1bp,y=-1bp,line cap=round,line join=round]
\path[use as bounding box] (0,0) rectangle (395,752);
\begin{pgfinterruptboundingbox}
\draw[line width=2.3bp] (75,215) .. controls (140,109) and (253,62) .. (326,56);
\draw[line width=1.8bp] (75,215) -- (327,218);
\draw[line width=1.7bp] (327,56) -- (327,85);
\draw[line width=1.7bp] (327,106) -- (326,257);
\draw[line width=1.7bp] (326,270) -- (324,316);
\draw[line width=1.7bp] (324,326) -- (323,411);
\draw[line width=1.7bp] (323,438) -- (317,684);
\node[inner sep=0pt,outer sep=0pt] at (315,37) {\scalebox{3.8}{\ensuremath{m}}};
\node[inner sep=0pt,outer sep=0pt] at (45,201) {\scalebox{3.8}{\ensuremath{M}}};
\node[inner sep=0pt,outer sep=0pt] at (347,207) {\scalebox{3.8}{\ensuremath{R}}};
\node[inner sep=0pt,outer sep=0pt] at (316,708) {\scalebox{3.8}{\ensuremath{C}}};
\end{pgfinterruptboundingbox}
\end{tikzpicture}
}\end{center}
Dieser Unterschied bei der Abweichung nach Osten — veranlasste mich, die Schlüsse von Dr.\ \textsc{Olbers} dazu genauer durchzugehen und nach der Ursache zu suchen. Sie liegt, wie mir scheint, darin, dass Dr.\ \textsc{Olbers} die wirkliche Bewegung des Körpers nach Osten \textit{allein} aus seiner anfänglichen Tangentialgeschwindigkeit ableitet und von der so entstehenden Bewegung die gleichzeitige Bewegung des Turmfußes abzieht, um die scheinbare Bewegung nach Osten zu erhalten. — Stellt aber die Fläche des Papiers die obige Ebene $\mathrm{A}$ dar, $\mathrm{C}$ den Mittelpunkt der Erde und $\mathrm{mM}$ die wirkliche Bewegung des Körpers, dann darf man meines Erachtens nicht außer Acht lassen, dass selbst die Anziehung nach $\mathrm{C}$ während der Bewegung nicht parallel zu $\mathrm{mC}$ ist. Eben des\SourceLeaf{503}halb wird die Geschwindigkeit nach Osten tatsächlich kleiner, sodass der Körper bei seiner Ankunft in $\mathrm{M}$ weniger weit nach Osten gelangt ist, als er mit der anfänglichen Geschwindigkeit gelangt wäre. Bei meiner Rechnung dazu finde ich auch, dass sich bei dieser Betrachtung die scheinbare Bewegung nach Osten tatsächlich um ein Drittel verringert.

\Needspace{7\baselineskip}
\begin{center}{\large\itshape Brief von Gauss an Benzenberg.}\end{center}
\begin{flushright}Braunschweig 1803. 8. März.\end{flushright}
—\quad—\quad—\quad—\quad Vor acht Tagen habe ich unserem Freund \textsc{Olbers} einen kleinen Aufsatz über die Abweichung fallender Körper geschickt. Heute erhalte ich darauf die Antwort:

1) Die Abweichung nach Osten sei nur $\tfrac23$ so groß wie die von ihm berechnete.

2) Meinen Schlüssen, dass die Abweichung nach Süden $=0$ sei, habe er nichts entgegenzusetzen. Er möchte aber wissen, \textit{worin} eigentlich \textit{sein} Gedankengang fehlerhaft sei.

Dazu bemerke ich noch Folgendes.

Vorausgesetzt, dass meine Schlüsse zur Abweichung nach Süden gewiss sind, scheint mir der Grund für die von Dr.\ \textsc{Olbers} gefundene Abweichung noch immer in seiner Voraussetzung zu liegen, der Körper \textit{bleibe} auch bei Luftwiderstand in der Ebene, die senkrecht zum Meridian steht und durch den Mittelpunkt der Erde geht. Diese Voraussetzung muss meines Erachtens unbedingt gerechtfertigt werden, aber ich bezweifle, dass sie sich rechtfertigen lässt. Durch die kegelförmige Bewegung der Luft gehen die Luftteilchen, in denen sich der Körper befindet, sobald die Erde ihre anfängliche Lage verlassen hat, in einem Winkel durch jene Ebene, den man nicht vernachlässigen darf. Dadurch tritt der Körper, dem die Luft nicht in Richtung dieser Ebene entgegenwirkt, nach Norden aus der Ebene heraus. Ich bin noch immer der Meinung, dass dadurch die Abweichung aus der Verzögerung vollständig ausgeglichen wird. Es scheint mir auch wahrscheinlich, dass \textsc{Guglielmini} genau dies sagen wollte und dass er nur deshalb bei \textsc{Olbers} keine Zustimmung fand, weil er sich nicht bestimmt genug ausdrückte. Ich hoffe jedoch zuversichtlich, dass entweder ich ganz zur Überzeugung von Dr.\ \textsc{Olbers} gelangen werde oder Dr.\ \textsc{Olbers} ganz zu meiner. —\quad—\quad—

\SourceLeaf{504}\Needspace{8\baselineskip}
\begin{center}{\large\itshape Fundamentalgleichungen für die Bewegung schwerer Körper\par auf der rotierenden Erde.}\end{center}
Die Lage eines Punktes wird auf zweierlei Weise bestimmt.

\textit{Erstens} durch seine senkrechten Abstände $X,Y,Z$ von drei zueinander senkrechten \textit{festen} Ebenen. Den gemeinsamen Schnittpunkt dieser Ebenen, $C$, legen wir in einen beliebigen Punkt der Erdachse. Die Ebene der $Z$ legen wir parallel zum Äquator, die Ebene der $Y$ in den Meridian, in dem der Anfangsort des Körpers liegt, und schließlich die Ebene der $X$ in den dazu senkrechten Meridian. Die $Z$ sind auf der Nordseite positiv, die $X$ auf der Seite des Anfangsorts des Körpers und die $Y$ auf der Seite, zu der dieser Anfangsort durch die Drehung geführt wird.

\textit{Zweitens} durch die senkrechten Abstände $x,y,z$ von drei zueinander senkrechten \textit{beweglichen} Ebenen, die also gegenüber der Erde ruhen und sich mit ihr drehen. Am zweckmäßigsten legen wir ihren gemeinsamen Schnittpunkt in den Anfangsort des Körpers. Die Ebene der $z$ legen wir senkrecht zur scheinbaren Richtung der Schwere, die der $y$ in den Meridian. Damit ist die zu beiden senkrechte Ebene der $x$ von selbst bestimmt. Die Pole dieser drei Ebenen sind also der Reihe nach der scheinbare Zenit, der Ostpunkt und der Südpunkt. \textit{Diese} Pole bezeichnen zugleich die Seiten der Ebenen, auf denen die Abstände $z,y,x$ positiv genommen werden.

Für den Punkt $C$ gelte nun $x=a$, $(y=0)$, $z=-c$. Ferner sei die (scheinbare, nördliche) Polhöhe des Beobachtungsorts $\varphi$ und der Winkel, um den sich die Erde nach der Zeit $t$ nach Osten bewegt hat, $\theta$. Unter diesen Voraussetzungen ergeben sich leicht die folgenden Gleichungen.
\[
\left.\begin{aligned}
x&=X\sin\varphi\cos\theta+Y\sin\varphi\sin\theta-Z\cos\varphi+a\\
y&=-X\sin\theta+Y\cos\theta\\
z&=X\cos\varphi\cos\theta+Y\cos\varphi\sin\theta+Z\sin\varphi-c
\end{aligned}\right\}\quad.\quad.\quad.\quad.\quad.\quad[1]
\]
\[
\left.\begin{aligned}
X&=(x-a)\sin\varphi\cos\theta-y\sin\theta+(z+c)\cos\varphi\cos\theta\\
Y&=(x-a)\sin\varphi\sin\theta+y\cos\theta+(z+c)\cos\varphi\sin\theta\\
Z&=-(x-a)\cos\varphi+(z+c)\sin\varphi
\end{aligned}\right\}\quad.\quad.\quad.\quad[2]
\]

\SourceLeaf{505}Die Koordinaten $X,Y,Z$ lassen sich einerseits als Funktionen von $t$ allein, andererseits auch als Funktionen der vier veränderlichen Größen $\theta,x,y,z$ auffassen. In der zweiten Betrachtungsweise haben sie daher vier partielle Differentiale. Es gilt also
\[
\begin{aligned}
\mathrm{d}X&=(\tfrac{\mathrm{d}X}{\mathrm{d}t})\mathrm{d}t=(\tfrac{\mathrm{d}X}{\mathrm{d}\theta})\mathrm{d}\theta+(\tfrac{\mathrm{d}X}{\mathrm{d}x})\mathrm{d}x+(\tfrac{\mathrm{d}X}{\mathrm{d}y})\mathrm{d}y+(\tfrac{\mathrm{d}X}{\mathrm{d}z})\mathrm{d}z\\
\mathrm{d}Y&=\mathrm{etc.}
\end{aligned}
\]
Die Geschwindigkeit des Körpers lässt sich ebenso wie seine Bewegung in drei Komponenten senkrecht zu den Ebenen der $X,Y,Z$ zerlegen. Diese sind demnach $(\tfrac{\mathrm{d}X}{\mathrm{d}t})$, $(\tfrac{\mathrm{d}Y}{\mathrm{d}t})$, $(\tfrac{\mathrm{d}Z}{\mathrm{d}t})$. Die Geschwindigkeiten des Luftelements hingegen, in dem er sich jeweils befindet, sind in Bezug auf dieselben Ebenen offenbar $(\tfrac{\mathrm{d}X}{\mathrm{d}\theta})\tfrac{\mathrm{d}\theta}{\mathrm{d}t}$, $(\tfrac{\mathrm{d}Y}{\mathrm{d}\theta})\tfrac{\mathrm{d}\theta}{\mathrm{d}t}$, $(\tfrac{\mathrm{d}Z}{\mathrm{d}\theta})\tfrac{\mathrm{d}\theta}{\mathrm{d}t}$. Die \textit{relativen} Geschwindigkeiten des Körpers in diesen drei Richtungen sind folglich
\[
\begin{aligned}
(\tfrac{\mathrm{d}X}{\mathrm{d}x})\tfrac{\mathrm{d}x}{\mathrm{d}t}+(\tfrac{\mathrm{d}X}{\mathrm{d}y})\tfrac{\mathrm{d}y}{\mathrm{d}t}+(\tfrac{\mathrm{d}X}{\mathrm{d}z})\tfrac{\mathrm{d}z}{\mathrm{d}t}&=\xi=\sin\varphi\cos\theta\tfrac{\mathrm{d}x}{\mathrm{d}t}-\sin\theta\tfrac{\mathrm{d}y}{\mathrm{d}t}+\cos\varphi\cos\theta\tfrac{\mathrm{d}z}{\mathrm{d}t}\\
(\tfrac{\mathrm{d}Y}{\mathrm{d}x})\tfrac{\mathrm{d}x}{\mathrm{d}t}+(\tfrac{\mathrm{d}Y}{\mathrm{d}y})\tfrac{\mathrm{d}y}{\mathrm{d}t}+(\tfrac{\mathrm{d}Y}{\mathrm{d}z})\tfrac{\mathrm{d}z}{\mathrm{d}t}&=\eta=\sin\varphi\sin\theta\tfrac{\mathrm{d}x}{\mathrm{d}t}+\cos\theta\tfrac{\mathrm{d}y}{\mathrm{d}t}+\cos\varphi\sin\theta\tfrac{\mathrm{d}z}{\mathrm{d}t}\\
(\tfrac{\mathrm{d}Z}{\mathrm{d}x})\tfrac{\mathrm{d}x}{\mathrm{d}t}+(\tfrac{\mathrm{d}Z}{\mathrm{d}y})\tfrac{\mathrm{d}y}{\mathrm{d}t}+(\tfrac{\mathrm{d}Z}{\mathrm{d}z})\tfrac{\mathrm{d}z}{\mathrm{d}t}&=\zeta=-\cos\varphi\tfrac{\mathrm{d}x}{\mathrm{d}t}+\sin\varphi\tfrac{\mathrm{d}z}{\mathrm{d}t}
\end{aligned}
\]
Die gesamte relative Geschwindigkeit ist folglich $=\BareRoot(\xi\xi+\eta\eta+\zeta\zeta)=u$. Wie man durch Entwicklung aus den obigen Werten leicht sieht, wird daraus $=\BareRoot(\tfrac{\mathrm{d}x^2}{\mathrm{d}t^2}+\tfrac{\mathrm{d}y^2}{\mathrm{d}t^2}+\tfrac{\mathrm{d}z^2}{\mathrm{d}t^2})$. Der Luftwiderstand ist ihrem Quadrat proportional. Wir setzen ihn daher $=Muu$ und zerlegen ihn nach den obigen drei Richtungen in $Mu\xi$, $Mu\eta$, $Mu\zeta$.

Wir betrachten hier die Erde als Rotationssphäroid. Die Richtung der Schwere geht daher durch die Erdachse. Der Punkt, an dem sie diese schneidet, liege um $q$ über $C$, sodass für ihn $Z=q$ gilt.

Setzt man nun außerdem die Stärke der Gravitation $=p$ und
\[
XX+YY+(Z-q)^2=rr
\]
so folgt aus den Grundsätzen der Dynamik
\SourceLeaf{506}\[
\left.\begin{aligned}
0&=\frac{\mathrm{dd}X}{\mathrm{d}t^2}+\frac{pX}{r}+Mu\xi\\
0&=\frac{\mathrm{dd}Y}{\mathrm{d}t^2}+\frac{pY}{r}+Mu\eta\\
0&=\frac{\mathrm{dd}Z}{\mathrm{d}t^2}+\frac{p(Z-q)}{r}+Mu\zeta
\end{aligned}\right\}\quad.\quad.\quad.\quad.\quad.\quad.\quad.\quad[3]
\]
Aus den obigen Werten von $X,Y,Z$ in [2] findet man, wenn man für $\tfrac{\mathrm{d}\theta}{\mathrm{d}t}$, das konstant ist, $n$ schreibt, die folgenden Gleichungen.
\[
\begin{aligned}
\frac{\mathrm{dd}X}{\mathrm{d}t^2}={}&\sin\varphi\cos\theta\frac{\mathrm{dd}x}{\mathrm{d}t^2}-\sin\theta\frac{\mathrm{dd}y}{\mathrm{d}t^2}+\cos\varphi\cos\theta\frac{\mathrm{dd}z}{\mathrm{d}t^2}\\
&-2n\sin\varphi\sin\theta\frac{\mathrm{d}x}{\mathrm{d}t}-2n\cos\theta\frac{\mathrm{d}y}{\mathrm{d}t}-2n\cos\varphi\sin\theta\frac{\mathrm{d}z}{\mathrm{d}t}-nnX\\[5pt]
\frac{\mathrm{dd}Y}{\mathrm{d}t^2}={}&\sin\varphi\sin\theta\frac{\mathrm{dd}x}{\mathrm{d}t^2}+\cos\theta\frac{\mathrm{dd}y}{\mathrm{d}t^2}+\cos\varphi\sin\theta\frac{\mathrm{dd}z}{\mathrm{d}t^2}\\
&+2n\sin\varphi\cos\theta\frac{\mathrm{d}x}{\mathrm{d}t}-2n\sin\theta\frac{\mathrm{d}y}{\mathrm{d}t}+2n\cos\varphi\cos\theta\frac{\mathrm{d}z}{\mathrm{d}t}-nnY\\[5pt]
\frac{\mathrm{dd}Z}{\mathrm{d}t^2}={}&-\cos\varphi\frac{\mathrm{dd}x}{\mathrm{d}t^2}+\sin\varphi\frac{\mathrm{dd}z}{\mathrm{d}t^2}
\end{aligned}\quad[4]
\]
Multipliziert man die drei Gleichungen [3] der Reihe nach mit $\sin\varphi\cos\theta$, $\sin\varphi\sin\theta$, $-\cos\varphi$ und addiert die Produkte, multipliziert man zweitens dieselben Gleichungen mit $-\sin\theta$, $\cos\theta$, $0$ und drittens mit $\cos\varphi\cos\theta$, $\cos\varphi\sin\theta$, $\sin\varphi$ und addiert beide Male die Produkte, so erhält man drei neue Gleichungen. Dabei setzt man für $\tfrac{\mathrm{dd}X}{\mathrm{d}t^2}$, $\tfrac{\mathrm{dd}Y}{\mathrm{d}t^2}$, $\tfrac{\mathrm{dd}Z}{\mathrm{d}t^2}$ ihre Werte aus [4] ein, für $X,Y,Z$ die Werte aus [2] und für $\xi$, $\eta$, $\zeta$ ihre jeweiligen Ausdrücke.
\[
\begin{aligned}
0&=\tfrac{\mathrm{dd}x}{\mathrm{d}t^2}-2n\sin\varphi\tfrac{\mathrm{d}y}{\mathrm{d}t}+(x-a)(\tfrac pr-nn)+\cos\varphi(\tfrac{pq}{r}-nnZ)+Mu\tfrac{\mathrm{d}x}{\mathrm{d}t}\\
0&=\tfrac{\mathrm{dd}y}{\mathrm{d}t^2}+2n\sin\varphi\tfrac{\mathrm{d}x}{\mathrm{d}t}+2n\cos\varphi\tfrac{\mathrm{d}z}{\mathrm{d}t}+y(\tfrac pr-nn)+Mu\tfrac{\mathrm{d}y}{\mathrm{d}t}\\
0&=\tfrac{\mathrm{dd}z}{\mathrm{d}t^2}-2n\cos\varphi\tfrac{\mathrm{d}y}{\mathrm{d}t}+(z+c)(\tfrac pr-nn)-\sin\varphi(\tfrac{pq}{r}-nnZ)+Mu\tfrac{\mathrm{d}z}{\mathrm{d}t}
\end{aligned}
\]
Ist der Körper also gegenüber der Erde in Ruhe, oder gilt $\mathrm{d}x=\mathrm{d}y=\mathrm{d}z=0$, so scheinen auf ihn senkrecht zu den Ebenen der $x,y,z$ die Kräfte
\SourceLeaf{507}\[
\begin{gathered}
(x-a)(\tfrac pr-nn)+\cos\varphi(\tfrac{pq}{r}-nnZ)\\
y(\tfrac pr-nn)\\
(z+c)(\tfrac pr-nn)-\sin\varphi(\tfrac{pq}{r}-nnZ)
\end{gathered}
\]
zu wirken. Auf einen bereits in Bewegung befindlichen Körper wirken die Kräfte dagegen anders. Der Luftwiderstand treibt den Körper in diesen Richtungen wie Kräfte der Stärke $Mu\tfrac{\mathrm{d}x}{\mathrm{d}t}$, $Mu\tfrac{\mathrm{d}y}{\mathrm{d}t}$, $Mu\tfrac{\mathrm{d}z}{\mathrm{d}t}$ und wirkt auf der rotierenden Erde somit genauso, wie er auf der ruhenden Erde wirken würde. Zu ihm kommen in jenen Richtungen noch die drei Kräfte
\[
-2n\sin\varphi\frac{\mathrm{d}y}{\mathrm{d}t},\quad 2n\sin\varphi\frac{\mathrm{d}x}{\mathrm{d}t}+2n\cos\varphi\frac{\mathrm{d}z}{\mathrm{d}t},\quad -2n\cos\varphi\frac{\mathrm{d}y}{\mathrm{d}t}
\]
hinzu, und allein durch diese wird die Rotation der Erde an fallenden Körpern sichtbar. Die bisherigen Schlüsse und Folgerungen sind streng und allgemein gültig.

Bei \textit{Versuchen}, die in dieser Hinsicht angestellt werden, bewegt sich der Körper stets in einem so kleinen Raum, dass man die Stärke der scheinbaren Schwere, die darin auf ruhende Körper wirkt, als unveränderlich $=g$ annehmen kann. Ihre Richtung kann man überall als parallel und damit als senkrecht zur Ebene der $z$ annehmen. Man darf daher ohne Bedenken anstelle der obigen drei Größen
\[
\begin{gathered}
(x-a)(\tfrac pr-nn)+\cos\varphi(\tfrac{pq}{r}-nnZ)\\
y(\tfrac pr-nn)\\
(z+c)(\tfrac pr-nn)-\sin\varphi(\tfrac{pq}{r}-nnZ)
\end{gathered}
\]
der Reihe nach $0,0,g$ einsetzen. Damit lauten die drei Fundamentalgleichungen
\[
\begin{aligned}
0&=\frac{\mathrm{dd}x}{\mathrm{d}t^2}-2n\sin\varphi\frac{\mathrm{d}y}{\mathrm{d}t}+Mu\frac{\mathrm{d}x}{\mathrm{d}t}\\
0&=\frac{\mathrm{dd}y}{\mathrm{d}t^2}+2n\sin\varphi\frac{\mathrm{d}x}{\mathrm{d}t}+2n\cos\varphi\frac{\mathrm{d}z}{\mathrm{d}t}+Mu\frac{\mathrm{d}y}{\mathrm{d}t}\\
0&=\frac{\mathrm{dd}z}{\mathrm{d}t^2}-2n\cos\varphi\frac{\mathrm{d}y}{\mathrm{d}t}+g+Mu\frac{\mathrm{d}z}{\mathrm{d}t}
\end{aligned}
\]
\SourceLeaf{508}Diese Gleichungen sind leicht zu integrieren, wenn man den Luftwiderstand vernachlässigt oder $M=0$ setzt. Man erhält
\[
\begin{aligned}
x&=\mathfrak{A}-\mathfrak{D}\cos\varphi.t+\tfrac{1}{2n}\mathfrak{E}\sin\varphi\cos(2nt+\mathfrak{F})+\tfrac12\sin\varphi\cos\varphi.gtt\\
y&=\mathfrak{B}-\tfrac{1}{2n}\mathfrak{E}\sin(2nt+\mathfrak{F})+\tfrac{1}{2n}\cos\varphi.gt\\
z&=\mathfrak{C}+\mathfrak{D}\sin\varphi.t+\tfrac{1}{2n}\mathfrak{E}\cos\varphi\cos(2nt+\mathfrak{F})-\tfrac12\sin\varphi^2.gtt
\end{aligned}
\]
Auch die folgenden Werte der willkürlichen Größen lassen sich leicht herleiten, wenn man voraussetzt, dass der Körper anfangs überhaupt keine scheinbare Geschwindigkeit hat.
\[
\begin{array}{lll}
\mathfrak{A}=-\tfrac{g}{4nn}\cos\varphi\sin\varphi,&\mathfrak{B}=0,&\mathfrak{C}=-\tfrac{g}{4nn}\cos\varphi^2\\
\mathfrak{D}=0,&\mathfrak{E}=\tfrac{g}{2n}\cos\varphi,&\mathfrak{F}=0
\end{array}
\]
Also
\[
\begin{aligned}
x&=\tfrac{g}{2n}\cos\varphi\sin\varphi(ntt-\tfrac{1}{2n}+\tfrac{1}{2n}\cos2nt)\\
y&=\tfrac{g}{2n}\cos\varphi(t-\tfrac{1}{2n}\sin2nt)\\
z&=-\tfrac12gtt+\tfrac{g}{2n}\cos\varphi^2(ntt-\tfrac{1}{2n}+\tfrac{1}{2n}\cos2nt)
\end{aligned}
\]
Diese Integration ist freilich nicht \textit{allgemein} zulässig, denn die obige Voraussetzung ist nur erlaubt, solange sich der Körper nicht weit von seinem anfänglichen scheinbaren Ort entfernt. In diesem Fall können wir aber die trigonometrischen Funktionen in Reihen entwickeln. Damit ergibt sich
\[
\begin{aligned}
x&=\tfrac16\cos\varphi\sin\varphi.gnnt^4\quad.\quad.\quad.\\
y&=\tfrac13\cos\varphi.gnt^3\quad.\quad.\quad.\\
z&=-\tfrac12gtt+\tfrac16\cos\varphi^2.gnnt^4\quad.\quad.\quad.
\end{aligned}
\]
Da die Fallzeit nur wenige Sekunden beträgt, also $nt$ höchstens einige Bogenminuten, und da (bei Radius $=1$) $1'=\tfrac{1}{3438}$ gilt, sind $x$ und der zweite Teil von $z$ ganz unmerklich. Also gilt $y=-\tfrac23z\cos\varphi.nt$. Bei Dr.\ \textsc{Benzenbergs} Versuch im Michaelisturm war $z=-235$ Fuß, $\varphi=53^0\,33'$ und $t=4''$ Sonnenzeit, also $nt=\tfrac{366}{365}$ Bogenminuten. Daraus ergibt sich $y=3{,}91$ Linien.

\SourceLeaf{509}Will man bei der Integration der obigen Gleichungen den Luftwiderstand berücksichtigen, so muss man sich mit Näherungen begnügen. Die Werte von $x,y,z$ lassen sich dann sehr leicht in Reihen nach den Potenzen von $n$ und $M$ entwickeln. Das höchste Glied von $x$ wird wie zuvor $=\tfrac16\cos\varphi\sin\varphi.gnnt^4$ und ist daher ganz ohne Bedeutung. Für $y$ und $z$ findet man unter Vernachlässigung der Quadrate und höheren Potenzen von $n$ und $M$ die folgenden Werte.
\[
\begin{aligned}
y&=\tfrac13\cos\varphi.gnt^3-\tfrac{1}{12}\cos\varphi.Mggnt^5\\
z&=-\tfrac12gtt+\tfrac{1}{12}Mggt^4.
\end{aligned}
\]
Setzen wir also $-z$, die wirkliche Fallstrecke, $=f$ und $\tfrac12gtt$, also die Fallstrecke im luftleeren Raum, $=f+\delta$, so gilt
\[
y=\tfrac23\cos\varphi.nt(f+\delta)-\cos\varphi.nt\delta=\tfrac23\cos\varphi.nt(f-\tfrac12\delta)
\]
Für die Versuche in St.\ Michael, wo $f+\delta=241{,}47$ Fuß war, erhalten wir daher die Abweichung nach Osten $y=3{,}86$ Linien.

